\documentclass[../../../include/open-logic-section]{subfiles}

\begin{document}

\olfileid{sth}{z}{infinity-again}	
\olsection{في اللانهاية}

تكفي المسلّمات السابقة لإثبات مجموعات غير متناهية العدد، بشرط وجود
مجموعة ما. فافرض وجود مجموعة؛ عندئذ توجد~$\emptyset$ بموجب
\olref[sth][z][sep]{prop:emptyexists}. ولكل مجموعة~$\OLId{latin-ordinary}{x}$ توجد
$\OLId{latin-ordinary}{x} \cup \{\OLId{latin-ordinary}{x}\}$ بموجب~\olref[sth][z][pairs]{prop:pairsconsequences}.
وبتكرار هذه العملية نحصل على:
\begin{enumerate}
	\item[\OLClassicalDigits{0}.] $\emptyset$ %& $0$ & stage $0$\\
	\item[\OLClassicalDigits{1}.] $ \{\emptyset\}$ %& $1$ & stage $1$\\
	\item[\OLClassicalDigits{2}.] $\{\emptyset, \{\emptyset\}\}$ %& $2$ & stage $2$\\
	\item[\OLClassicalDigits{3}.] $\{\emptyset, \{\emptyset\}, \{\emptyset, \{\emptyset\}\}\}$ %&$3$ & stage $3$\\
	\item[\OLClassicalDigits{4}.] $\{\emptyset, \{\emptyset\}, \{\emptyset, \{\emptyset\}\}, \{\emptyset, \{\emptyset\}, \{\emptyset, \{\emptyset\}\}\}\}$% & $4$ & stage $4$ 
\end{enumerate}
ويسهل التحقق من تباين هذه المجموعات.

وقد ابتدأنا العدد بـ~$\OLMathNumeral{0}$ لوجوه؛ منها سهولة التحقق من أن المجموعة
المسماة «$\OLId{latin-ordinary}{n}$» ذات~$\OLId{latin-ordinary}{n}$ من العناصر سواء، وأنها، بحسب التصور، تتكون في
المرحلة رقم~$\OLId{latin-ordinary}{n}$.

غير أن ذلك يثبت \emph{مجموعات غير متناهية العدد}، ولا يثبت
\emph{مجموعة غير منتهية} كثيرة العناصر. والفرق جوهري؛ فإن لم نجد
مجموعة لا نهائية بمعنى ددكند تعذر بناء جبر ددكندي، ونحن نحتاج إلى هذا
الجبر لنتخذه مجموعة الأعداد الطبيعية. قارن
\olref[sfr][infinite][dedekindsproof]{sec}.

والأهم أن المسلّمات الموضوعة إلى هنا \emph{لا} تكفل مجموعة لا نهائية؛
فلا بد من مسلّمة جديدة:

\begin{axiom}[اللانهاية]
توجد مجموعة~$\OLId{latin-looped}{I}$ تحقق~$\emptyset \in \OLId{latin-looped}{I}$، ويكون
$\OLId{latin-ordinary}{x} \cup \{\OLId{latin-ordinary}{x}\} \in \OLId{latin-looped}{I}$ كلما كان~$\OLId{latin-ordinary}{x} \in \OLId{latin-looped}{I}$.
\begin{align*}
	\exists \OLId{latin-looped}{I}( & (\exists \OLId{latin-ordinary}{o} \in \OLId{latin-looped}{I})\forall \OLId{latin-ordinary}{x}\ \OLId{latin-ordinary}{x} \notin \OLId{latin-ordinary}{o} \land {}\\
	& (\forall \OLId{latin-ordinary}{x} \in \OLId{latin-looped}{I})(\exists \OLId{latin-ordinary}{s} \in \OLId{latin-looped}{I})\forall \OLId{latin-ordinary}{z}(\OLId{latin-ordinary}{z} \in \OLId{latin-ordinary}{s} \liff (\OLId{latin-ordinary}{z} \in \OLId{latin-ordinary}{x} \lor \OLId{latin-ordinary}{z} = \OLId{latin-ordinary}{x})))
\end{align*}
\end{axiom}

وظاهر أن المجموعة~$\OLId{latin-looped}{I}$ التي تقررها مسلّمة اللانهاية لا نهائية بمعنى
ددكند: عنصرها المميز~$\emptyset$، والتطبيق الحقني على
$\OLId{latin-looped}{I}$ معطى بـ$\OLId{latin-ordinary}{s}(\OLId{latin-ordinary}{x}) = \OLId{latin-ordinary}{x}\cup\{\OLId{latin-ordinary}{x}\}$. وقد شرحت
\olref[sfr][infinite][dedekind]{thm:DedekindInfiniteAlgebra} استخراج
جبر ددكندي من مجموعة لا نهائية بمعنى ددكند؛ وهذا الجبر هو الذي سنعامله
معاملة مجموعة الأعداد الطبيعية. وعلى الضبط:
\begin{defn}\ollabel{defnomega}
لتكن~$\OLId{latin-looped}{I}$ مجموعة تقررها مسلّمة اللانهاية، ولتكن~$\OLId{latin-ordinary}{s}$ الدالة
$\OLId{latin-ordinary}{s}(\OLId{latin-ordinary}{x}) = \OLId{latin-ordinary}{x} \cup \{\OLId{latin-ordinary}{x}\}$. ولتكن~$\omega = \closureofunder{\OLId{latin-ordinary}{s}}{\emptyset}$.
نسمي عناصر~$\omega$ \emph{الأعداد الطبيعية}، ونقول إن~$\OLId{latin-ordinary}{n}$ حاصل تطبيق
$\OLId{latin-ordinary}{s}$ على~$\emptyset$ عدد~$\OLId{latin-ordinary}{n}$ من المرات.
\end{defn}

وبالرجوع إلى القائمة السابقة يتبين أن المجموعة الموسومة «$\OLId{latin-ordinary}{n}$» تُعامل
\emph{بوصفها} العدد~$\OLId{latin-ordinary}{n}$.

وسيأتي الكلام على شأن هذا الاصطلاح في~\olref[sth][z][nat]{sec}. وهو
يمكننا الآن من إثبات نتيجة ظاهرة:

\begin{prop}\ollabel{naturalnumbersarentinfinite}
لا يكون عدد طبيعي غير متناهٍ على معنى ددكند.
\end{prop}

\begin{proof}
نسلك الاستقراء المقرر في
\olref[sfr][infinite][induction]{thm:dedinfiniteinduction}. ومن البيّن
أن~$\OLMathNumeral{0} = \emptyset$ ليس لا نهائيًا بمعنى ددكند. وفي خطوة الاستقراء نثبت
النقيض المقابل: لو كان~$\OLId{latin-ordinary}{s}(\OLId{latin-ordinary}{n})$، على المحال، لا نهائيًا بمعنى ددكند،
لكان~$\OLId{latin-ordinary}{n}$ كذلك.

فافرض~$\OLId{latin-ordinary}{s}(\OLId{latin-ordinary}{n})$ لا نهائيًا بمعنى ددكند، أي يوجد !!{injection}~$\OLId{latin-ordinary}{f}$
بحيث~$\ran{\OLId{latin-ordinary}{f}}\subsetneq \dom{\OLId{latin-ordinary}{f}} = \OLId{latin-ordinary}{s}(\OLId{latin-ordinary}{n}) = \OLId{latin-ordinary}{n} \cup \{\OLId{latin-ordinary}{n}\}$. وهناك حالتان
ينبغي النظر فيهما.

\emph{الحالة \OLClassicalDigits{1}: $\OLId{latin-ordinary}{n} \notin \ran{\OLId{latin-ordinary}{f}}$.} حينئذ~$\ran{\OLId{latin-ordinary}{f}} \subseteq \OLId{latin-ordinary}{n}$
و$\OLId{latin-ordinary}{f}(\OLId{latin-ordinary}{n}) \in \OLId{latin-ordinary}{n}$. ضع~$\OLId{latin-ordinary}{g} = \funrestrictionto{\OLId{latin-ordinary}{f}}{\OLId{latin-ordinary}{n}}$؛ فيكون
$\ran{\OLId{latin-ordinary}{g}} = \ran{\OLId{latin-ordinary}{f}} \setminus \{\OLId{latin-ordinary}{f}(\OLId{latin-ordinary}{n})\} \subsetneq \OLId{latin-ordinary}{n} = \dom{\OLId{latin-ordinary}{g}}$.
ومن ثم يكون~$\OLId{latin-ordinary}{n}$ لا نهائيًا بمعنى ددكند.

\emph{الحالة \OLClassicalDigits{2}: $\OLId{latin-ordinary}{n} \in \ran{\OLId{latin-ordinary}{f}}$.} اختر
$\OLId{latin-ordinary}{m} \in \dom{\OLId{latin-ordinary}{f}} \setminus \ran{\OLId{latin-ordinary}{f}}$، وعرّف دالة~$\OLId{latin-ordinary}{h}$ مجالها
$\OLId{latin-ordinary}{s}(\OLId{latin-ordinary}{n}) = \OLId{latin-ordinary}{n} \cup \{\OLId{latin-ordinary}{n}\}$:
\[
\OLId{latin-ordinary}{h}(\OLId{latin-ordinary}{x}) = 
	\begin{cases}
		\OLId{latin-ordinary}{f}(\OLId{latin-ordinary}{x}) & \text{إذا كان }\OLId{latin-ordinary}{f}(\OLId{latin-ordinary}{x}) \neq \OLId{latin-ordinary}{n}\\
		\OLId{latin-ordinary}{m} & \text{إذا كان }\OLId{latin-ordinary}{f}(\OLId{latin-ordinary}{x})=\OLId{latin-ordinary}{n}
	\end{cases}
\]
فتتفق~$\OLId{latin-ordinary}{h}$ و$\OLId{latin-ordinary}{f}$ في كل موضع إلا في أن
$\OLId{latin-ordinary}{h}(\OLId{latin-ordinary}{f}^{-\OLMathNumeral{1}}(\OLId{latin-ordinary}{n})) = \OLId{latin-ordinary}{m} \neq \OLId{latin-ordinary}{n} = \OLId{latin-ordinary}{f}(\OLId{latin-ordinary}{f}^{-\OLMathNumeral{1}}(\OLId{latin-ordinary}{n}))$. ولما كانت~$\OLId{latin-ordinary}{f}$
!!a{injection}، لزم~$\OLId{latin-ordinary}{n} \notin \ran{\OLId{latin-ordinary}{h}}$، وكذلك
$\ran{\OLId{latin-ordinary}{h}} \subsetneq \dom{\OLId{latin-ordinary}{h}} = \OLId{latin-ordinary}{s}(\OLId{latin-ordinary}{n})$. فحجة الحالة~\OLClassicalDigits{1} تجعل~$\OLId{latin-ordinary}{n}$
لا نهائيًا بمعنى ددكند.
\end{proof}

ويبقى بعد ذلك مسوغ مسلّمة اللانهاية. والجواب المختصر أن روايتنا تحتاج
إلى أصل آخر، هو:
\begin{enumerate}
	\item[] \stagesinf. توجد مرحلة لا نهائية. أي توجد مرحلة (أ) ليست المرحلة الأولى، و(ب) تسبقها بعض المراحل، لكنها (ج) ليس لها سابق مباشر.
\end{enumerate}
فتلزم مسلّمة اللانهاية عن هذا الأصل مباشرة؛ لأن العدد الطبيعي~$\OLId{latin-ordinary}{n}$
يتكون في المرحلة~$\OLId{latin-ordinary}{n}$، فتتكون المجموعة~$\omega$ عند أول مرحلة لا
نهائية، وتكون~$\omega$ شاهدًا للمسلّمة.

لكن السؤال يعود في مسوغ~\stagesinf. فهو، مثل~\stagessucc، غير منصوص
عليه في رواية الهرم التراكمي-التكراري. بل لا شيء في مجرد تكون
المجموعات مرحلة بعد مرحلة يضطرنا إلى إثبات مرحلة \emph{لا نهائية}؛
إذ لا تناقض في تصور المراحل مرتبة كالأعداد الطبيعية.

وهنا تظهر معضلة. فقد بحثنا في
\olref[sfr][infinite][dedekindsproof]{sec} «برهان» ددكند على وجود
مجموعة من الأفكار لا نهائية بمعنى ددكند، ولعله لم يقنعك. فإذا لم يكن
\stagesinf{} لازمًا عن التصور التكراري للمجموعة، أو عن «قوانين الفكر»،
بقينا من غير مسوغ داخلي لوجود مجموعة لا نهائية بمعنى ددكند.

وفي المسألة كلام كثير وراء هذا. وحسبنا أن يكون القارئ قد بلغ موضعًا
يستطيع فيه النظر فيما \emph{يلزم} لتسويغ مسلّمة أو مبدأ. وسنسلم فيما
يأتي بـ~\stagesinf{}.

\end{document}
